\documentclass[10pt,a4paper]{amsart}
\usepackage[margin=19mm]{geometry}
\usepackage{amsmath,amssymb,mathtools}
\usepackage[scr=rsfs,cal=euler]{mathalfa}
\usepackage{xfrac}
\usepackage[unicode,hidelinks,bookmarks=false]{hyperref}
\newcommand{\ie}{i.e.\ }
\newcommand{\cf}{cf.\ }
\newcommand{\resp}{respectively\ }
\newcommand{\Subsection}{\subsection}
\providecommand{\nobreakdash}{\nobreak\hbox{-}}
\setlength{\emergencystretch}{2em}
\multlinegap0pt
\usepackage{amssymb}
\usepackage{float}
\usepackage{bm}
\usepackage{mathtools}
\usepackage[shortlabels]{enumitem}
\newtheorem{lemma}{Lemma}
\newtheorem{theorem}{Theorem}
\renewcommand{\Im}{\operatorname{Im}}
\renewcommand{\Re}{\operatorname{Re}}
\def \br{\mathbf{r}}
\def \inv{^{-1}}
\title{Efficient time-domain scattering synthesis via\\ frequency-domain singularity subtraction}
\author{Oscar P. Bruno, Manuel A. Santana}
\date{Source: arXiv:2505.06189v2 (CC BY 4.0)}
\begin{document}
\maketitle

\noindent\emph{Source and licence.} Efficient time-domain scattering synthesis via\\ frequency-domain singularity subtraction. Version arXiv:2505.06189v2 (CC BY 4.0), distributed by arXiv under the Creative Commons Attribution 4.0 International licence, http://creativecommons.org/licenses/by/4.0/, as recorded in the arXiv licence metadata for this exact version and shown on the abstract page https://arxiv.org/abs/2505.06189v2. Full licence notice: \texttt{licenses/arxiv-2505.06189-CC-BY-4.0.txt}.\par\smallskip

\noindent\textit{Editorial scope.} Counted unit: the article's theorem thm:app (source0.tex lines 2826--2832). The article's complete original proof of it is delivered with it (source0.tex lines 2833--2861, one \texttt{proof} environment of 29 lines). No mathematics is written, repaired or re-derived: the statement and the proof are the article's own lines, in the article's own notation, and no retained line is substituted or re-flowed.  Retained uncounted as the source-local context the proof needs: lines 551--553; lines 847--849; lines 2798--2803.  Nothing was omitted from any retained span, so the package carries no omission record.  The retained lines cite 1 works, whose entries are transcribed verbatim from the article's own bibliography source in the e-print and delivered with short editorial labels recorded in the item's reference mapping.  No retained line contains a figure environment, an image-inclusion command or drawing code, so no image is delivered and the package declares no asset dependency.  Editorial formatting only, in addition: an AMS \texttt{amsart} preamble that restates the article's own declarations and macros used by the retained lines, namely the environments \texttt{lemma}, \texttt{theorem} on separate counters, together with the standard packages those lines need.  The retained environments print in this document as Lemma 1, Theorem 1 in order of appearance, whereas the article prints the same items with the numbering of its own full document.  The complete native source of the article as distributed in the arXiv e-print for this exact version is bundled unchanged as \texttt{sources/177871/source0.tex}.
 \begin{equation}\label{eqn:sl_dl_operator}
    (S_\omega\psi)(\br) = \int_\Gamma G_\omega(\br,\br')\psi(\br',\omega)d\sigma(r') \quad \text{and} \quad (K_\omega \psi)(\br) = \int_{\Gamma}\frac{\partial G_\omega(\br,\br')}{\partial n(\br')}\psi(\br',\omega)d\sigma(r'), \quad \br \in \Gamma,
\end{equation}

\begin{equation}\label{eqn:sol_op_comb}
  \mathcal{U}^\mathrm{c}_\omega[B] =  \mathcal{K}_\omega[(C_{\omega,\eta})\inv B](\omega) - \mathrm{i}\eta \mathcal{S}_\omega[(C_{\omega,\eta})\inv B].
\end{equation}

\begin{lemma}\label{lem:Injectivty}
  Let $\omega$ and $\eta$ satisfy $\Im(\omega) < 0$, $\Re(\omega) > 0$
  (resp. $\Re(\omega) < 0$), and $\eta < 0$ (resp. $\eta > 0$)
  . Then $\mathcal{C}_\eta : H^{1/2}(\Gamma) \to H^1_{loc}(\Omega^e)$
  is an injective operator.
\end{lemma}

\begin{theorem}\label{thm:app}
  Let $\omega \in \mathbb{C}$ such that $\Re(\omega) > 0$
  (resp. $\Re(\omega) < 0$), and let $\eta < 0$ (resp. $\eta >
  0$). Then, for $\Im(\omega) \leq 0$, the set of poles of
  $(C_{\omega,\eta})^{-1}$ coincides with the set of poles of
  $\mathcal{U}_{\omega}^\mathrm{c}$.
\end{theorem}

\begin{proof}
  Since the double- and single-layer
  operators~\eqref{eqn:sl_dl_operator} are compact, the operator
  $(C_{\omega,\eta})^{-1}$ is a meromorphic function of $\omega$ in
  the entire complex plane~\cite[Proposition 7.4]{taylorbook}, except
  for a logarithmic branch cut joining $\omega = 0$ and
  $\omega =\infty$. In view of the
  representation~\eqref{eqn:sol_op_comb} of the solution operator
  $\mathcal{U}^\mathrm{c}_\omega$ it follows that the set of poles of
  $\mathcal{U}_\mathrm{c}(\omega)$ is contained within the set of
  poles of $(C_{\omega,\eta})^{-1}$.

  To show that the converse is also true assume
  $(C_{\omega,\eta})\inv$ has a pole of order $m$ at
  $\omega = \omega_0$ with $\Im(\omega_0) < 0$.  Then there exists an
  element $B \in H^{1/2}(\Gamma)$ such that
    \begin{equation*}
      (C_{\omega,\eta})\inv [B] = (\omega - \omega_0)^{-m}\left(B_m + B_{m+1}(\omega - \omega_0) + \cdots\right)
      \end{equation*}
      for a certain sequence $B_j\in H^{1/2}(\Gamma)$, $j\geq m$, with
      $B_m\ne 0.$ Letting $u_m = \mathcal{C}_\eta[B_m]$ it follows
      that
     \begin{equation*}
        \mathcal{U}^\mathrm{c}_\omega [B] = (\omega - \omega_0)^{-m}u_m + O((\omega - \omega_0)^{-m +1}) \quad \mbox{as}\quad \omega \to \omega_0.
     \end{equation*}
     By Lemma~\ref{lem:Injectivty} $u_m \ne 0$, and, therefore
     $\omega_0$ is a pole of $\mathcal{U}^\mathrm{c}_\omega$. The
     proof is now complete.
\end{proof}

\begin{thebibliography}{99}
\bibitem[taylor]{taylorbook} Michael~E. Taylor. \newblock {\em Partial differential equations. {II}}. \newblock Springer-Verlag, New York, 1996.
\end{thebibliography}
\end{document}
